% Use cases to test the revised (two-tier) definition of readings and reading sequences against. Scenarios only, no proposed fixes.
\documentclass{amsart}
\usepackage{hyperref}
\usepackage{amssymb}
\usepackage{enumitem}
\newlist{todo}{itemize}{2}
\setlist[todo]{label=$\square$}
\title{Reading Sequences: Use Cases to Think About}
\begin{document}
\maketitle
\begin{abstract}
An itinerary of scenarios, mostly involving mutable data, against which to check the two-tier definition of readings (tier~1, over a container) and reading sequences (tier~2, sequences of container--reading pairs), together with the stratification rule that a container's canonical reading must be a tier-1 reading over that container.
\end{abstract}

\section{Promoting a reading to canonical}
\begin{todo}
\item Alice wants to make Bob's tier-2 sequence of her series her canonical reading. Bob's sequence uses a helper container $X=\{DT,SC,MN\}$ that he minted for an ``any order'' step.
\item Alice tries to re-express Bob's sequence as a tier-1 reading. To get ``any order'' over three of her children, she needs $X$, or her own copy of it, among her children.
\item Alice adds $X$ as a child of her container. Those three books are now in it twice, once directly and once through $X$. What happens to her natural and trivial readings, and to the canonical containment chains of those books?
\item Alice wants an interleaved canonical reading: book~1 ch.~1--5, then book~2, then book~1 ch.~6--10.
\end{todo}

\section{Adding and removing children}
\begin{todo}
\item Alice has a custom canonical reading and posts chapter~12. The reading doesn't mention the new chapter.
\item Alice adds a child and updates the canonical reading in two separate writes. What do readers see between the two writes, or if she forgets the second?
\item Same as above, but the new child lives in someone else's space repo.
\item Alice removes chapter~7. Her canonical reading still references $c_7$, and so do tier-1 readings over $A$ written by other people.
\item Alice's canonical reading becomes invalid. What do readers get instead?
\item Alice reorders her container's children for display reasons. Every sequence that cites ``natural reading of $A$'' changes.
\end{todo}

\section{Edits that disqualify a canonical reading}
\begin{todo}
\item Alice edits her canonical reading of $A$ to add a step over a container that isn't one of $A$'s children, such as a sibling or Bob's helper container. It may still be a reading of $A$ by atoms, but it is now shaped like a tier-2 sequence and can't be canonical.
\item Alice edits her canonical reading of $A$ to add the step ``read $A$ canonically''.
\item One reading record is $A$'s canonical reading and is also cited as a step in other people's tier-2 sequences. An edit that is fine for those tier-2 uses disqualifies it as canonical. Does the edit break one use and not the other?
\item Alice designates a tier-1 reading Bob wrote over $A$ as her canonical reading. Bob later edits it so it no longer qualifies.
\item A reading is re-pointed or re-created against another container record that describes the same set as $A$. It is no longer a reading over $A$.
\item An edit leaves the canonical reading valid but changes its character: it becomes infinite, or stops reading some children.
\end{todo}

\section{Attribution moves without anyone touching the reading}
\begin{todo}
\item Alice moves chapter~7 into a ``Deleted scenes'' container $D$ outside $B$. Fan sequences that touch ch.~7 are no longer readings of $B$.
\item Alice folds a standalone one-shot into $B$ as a bonus chapter. Where do existing readings of the one-shot now surface?
\item The reverse: Alice pulls a chapter out of $B$ and publishes it as a standalone. Where do its readings surface?
\item Alice splits chapter~3 into 3a and 3b, deleting the old atom. Every reference to $\{c_3\}$ dangles.
\item Alice merges two chapters into one atom.
\end{todo}

\section{Revisions and moves}
\begin{todo}
\item Alice revises a chapter: a typo fix, or a substantial rewrite.
\item Bob's tier-2 sequence reads Alice's ch.~3, then his missing scene. Alice rewrites ch.~3 so the missing scene no longer makes sense.
\item Alice drafts a fic in a private beta-reader Space and builds its canonical reading there. She then publishes it to the public archive Space.\footnote{A record's URI includes its Space: \nolinkurl{at://{spaceDid}/space/{spaceType}/{skey}/{authorDid}/{collection}/{rkey}}. See \url{https://github.com/bluesky-social/proposals/tree/main/0016-permissioned-data}.}
\end{todo}

\section{Changes made by other people}
\begin{todo}
\item Alice's anthology reading reads Bob's story by its canonical reading. Bob adds an epilogue Alice never curated.
\item Bob changes his story's canonical reading to an infinite terminal loop. The steps after it in Alice's anthology reading, and in tier-2 sequences that read the anthology, can no longer be reached.
\item Alice adds $B$ to $A$ while Bob, at the same time, adds $A$ to $B$. Each edit is fine on its own, but together they make a containment cycle.
\item Two app-views see those two edits in different orders.
\end{todo}

\section{Deletion, moderation and access}
\begin{todo}
\item Bob deletes his story from Alice's anthology. Alice isn't notified, and readings over the anthology break.
\item Bob's story is taken down by moderation. Alice's whole anthology, and every sequence touching it, is hidden.
\item A reader blocks Bob. Is Alice's anthology hidden for that reader?
\item Bob is removed from the Space's writer set.
\item Bob deletes or blocks on purpose to grief Alice's anthology.
\item Bob's sequence uses Alice's alternate ``chronological'' reading $R$, and Alice deletes $R$.
\end{todo}

\section{Reader state}
\begin{todo}
\item A reader is at step $n$ of a reading, and a step is inserted before $n$.
\item Alice changes her canonical reading while a reader is halfway through it. What does ``next chapter'' mean?
\item ``Next unread'' in a reading where atoms repeat, as in \emph{Rayuela} or loops.
\item A reader is partway through a trivial (``any order'') reading. The app must remember the order they actually took.
\item That reader's trivial reading gains or loses children mid-read.
\end{todo}

\end{document}
